\subsection{The \alert{JSON Structure} of the \texttt{<\alert{mapping-section}>}s}

\begin{frame}[fragile]{The \alert{JSON Structure} of the \texttt{<\alert{mapping-section}>}s}
  A <\texttt{\textcolor{fh-ooe-red}{mapping-section}}> object can be one of the three objects \texttt{\textcolor{fh-ooe-red}{static}}, \texttt{\textcolor{fh-ooe-red}{value}}, and \texttt{\textcolor{fh-ooe-red}{json}}.
  \begin{itemize}
    \item Only one of these objects needs to be specified.
    \item If two or all three objects are specified, the order of interpretation is
    \begin{enumerate}
      \item \texttt{\textcolor{fh-ooe-red}{static}}
      \item \texttt{\textcolor{fh-ooe-red}{value}}
      \item \texttt{\textcolor{fh-ooe-red}{json}}
    \end{enumerate}
    \item \alert{All three object describe a topic/message mapping} but in different ways.
    \begin{description}[\texttt{static}]
      \item[\texttt{static}] is thought for a \alert{string-matching} mapping of incoming messages to outgoing one.
      \item[\texttt{value}] is used if the incoming message should be \alert{interpreted as a value} and accessible. The actual mapping is done controlled by a \alert{\href{https://github.com/pantor/inja}{inja}} template, whereas the message is accessible via the template variable \texttt{\alert{message}}.
      \item[\texttt{json}] can be used in case the \alert{incoming message} is a \alert{JSON object}. The \alert{actual mapping} is done controlled by a \alert{\href{https://github.com/pantor/inja}{inja}} template. The JSON objects are accessible in the template \alert{as elements of} the template JSON variable \texttt{\alert{message}}.
    \end{description}
  \end{itemize}
\end{frame}

